\subsection{Joint traffic and execution measurements}\label{sec:cover-study}
The study compares six cover policies on matched synthetic support tickets. Two fixed-weight models, Qwen3:8B and Llama3.1:8B, each run 300 fixtures per condition: 120 training, 60 validation and 120 test, balanced across six actions. The 3,600 episodes have terminal results. Fixtures and cover draws are paired across conditions; request phrasings recur across splits.

\paragraph{Execution and observation}
The models choose checks, business actions, objects, recovery and termination through LangGraph~\cite{langchain2025langgraph}. A controller fetches the ticket and enforces prerequisites without supplying the correct action. Seven local MCP Streamable HTTP destinations use TLS~1.3/HTTP/2, a common \path{/mcp} path and request/response padding quanta of 4,096/16,384 bytes. The observer receives destinations and bidirectional ciphertext byte counts within supplied episode and call windows. Timing and model traffic are outside this view.

A frozen schedule inserts one or two random decoys among early tool calls. The conditions are no cover, proxy discard, read cover, isolated execution, live write, and isolated execution with visible responses. Proxy discard computes on a disposable copy; isolated execution retains a separate clone. Live writes deliberately modify the primary fixture, including other orders or shared quota. The visible control sends clearly marked cover responses from the clone to the planner. Other cover responses are hidden. The live-write arm is an adverse control selected to exercise collateral changes. Proxy discard also computes a disposable response; it is not a pure transport-only dummy.

\paragraph{Matched endpoints and attacks}
\emph{Goal completion} checks the requested action/object, prerequisites and target postcondition. \emph{State integrity} checks the entire business state against the single intended-action reference. We also count completed tasks with state errors directly; the two predicates are not nested. Recognition predicts the requested action including failed executions. Logistic regression (LR) and random forest (RF) train within each condition, with validation selecting the primary attacker. Frozen and destination-position attacks are retained alongside them.

Each endpoint has a prespecified twelve-test Holm family: five cover-minus-baseline contrasts and visible-minus-hidden for each model. Conditional intervals resample paired fixtures by action, holding fitted attacks fixed. Supplement S8 contains every attack and contrast, state replay and a worked evaluation protocol.

\begin{table*}[t]\centering\small
\caption{Executed cover: all held-out conditions, 120 paired fixtures per model and condition. Recognition columns count correct requested-action predictions; chance is 20/120. The learner is selected using validation LR/RF accuracy. The position rule remains visible when it is stronger. Goal, whole-state integrity and completed tasks with state errors are separate endpoints.}
\label{tab:cover-main}
\setlength{\tabcolsep}{4pt}
\begin{tabularx}{\textwidth}{@{}lYrrrrrr@{}}\toprule
Model & Cover policy & Selected & Last dest. & Goal & State & Goal + error & Call bytes $\times$ \\\midrule
Qwen3:8B & No cover & 118 & 116 & 114 & 116 & 0 & 1.00 \\
Qwen3:8B & Proxy discard & 62 & 77 & 114 & 116 & 0 & 1.34 \\
Qwen3:8B & Read cover & 62 & 77 & 114 & 115 & 0 & 1.34 \\
Qwen3:8B & Isolated, hidden & 61 & 77 & 114 & 115 & 0 & 1.34 \\
Qwen3:8B & Live write & 62 & 75 & 102 & 7 & 95 & 1.34 \\
Qwen3:8B & Isolated, visible & 63 & 77 & 113 & 114 & 1 & 1.35 \\
\midrule
Llama3.1:8B & No cover & 113 & 79 & 82 & 69 & 13 & 1.00 \\
Llama3.1:8B & Proxy discard & 79 & 67 & 80 & 68 & 13 & 1.29 \\
Llama3.1:8B & Read cover & 72 & 67 & 84 & 71 & 14 & 1.28 \\
Llama3.1:8B & Isolated, hidden & 82 & 67 & 83 & 70 & 13 & 1.28 \\
Llama3.1:8B & Live write & 80 & 65 & 76 & 3 & 73 & 1.29 \\
Llama3.1:8B & Isolated, visible & 81 & 68 & 83 & 73 & 11 & 1.29 \\
\bottomrule\end{tabularx}
\end{table*}


\paragraph{Recognition follows the schedule, not the cover semantics}
Every cover policy reduces requested-action recognition, and all ten primary cover-versus-baseline contrasts survive correction. The four hidden-response policies differ only in what the decoy executes. For Qwen they leave the selected attacker at 62, 62, 61 and 62 of 120 (proxy discard, read, isolated, live write); for Llama the range is 72--82. Executing deeper therefore gave no visible privacy benefit to this observer. These are point estimates, and we do not claim equivalence. What does limit protection is coverage. The last-destination rule reaches 77/120 for Qwen under every hidden policy, above the selected learner, because the frozen schedule places decoys only among early calls. A stronger defense under this observer extends schedule coverage to the end of the trace; making each decoy more realistic does not help.

\paragraph{State integrity follows the cover semantics}
The same policies differ sharply in what they do to the application. Proxy discard, read and isolated cover keep Qwen's whole-state integrity at 115--116/120, close to no cover. Live writes reduce it to 7/120, and for Llama from 69/120 to 3/120. Both reductions survive correction. Completion falls by only 10.0 and 5.0 points, and only Qwen's drop is significant. Among completed tasks, 95 of 102 (Qwen) and 73 of 76 (Llama) carry collateral state errors, against 0 of 114 and 13 of 82 without cover. A completion-only utility metric would therefore have scored the most damaging policy as nearly harmless. We chose live writes as an adverse control in order to exercise this gap. The counts show what completion misses; they do not estimate how often deployed cover causes harm.

\paragraph{Design consequence and cost}
Together, the two contrasts favor the cheapest boundary that the observer model permits. For Qwen, proxy discard matches isolated execution in recognition (62 vs.\ 61), completion (114/120) and call bytes (1.34$\times$), and it retains no decoy state. Isolation or execution is justified only when the observer can see something that requires realistic downstream work, such as response timing or third-party calls. That case lies outside this tested view.

Making isolated cover responses visible to the planner raises mean prompt tokens by 18.1\% and 18.9\%. All primary visible-minus-hidden contrasts have adjusted $p=1$. Planner paths vary between paired runs, and we have no repeated-condition cohort, so this average effect remains unresolved (Supplement S8).
